\documentclass[11pt,a4paper]{article}
\usepackage[utf8]{inputenc}
\usepackage[vietnamese]{babel}
\usepackage{amsmath}
\usepackage{amssymb}
\usepackage{geometry}
\usepackage{booktabs}
\usepackage{xcolor}
\usepackage{hyperref}
\usepackage{enumitem}

\geometry{top=2.5cm, bottom=2.5cm, left=2.5cm, right=2.5cm}

\title{\textbf{HƯỚNG DẪN ÔN TẬP BÀI GIẢNG VÀ PHƯƠNG PHÁP GIẢI\\KINH TẾ HỌC VI MÔ (CHƯƠNG 1 - 4)}}
\author{Hệ Thống Trắc Nghiệm Kinh Tế Học Vi Mô}
\date{Tháng 6, Năm 2026}

\begin{document}

\maketitle
\tableofcontents
\newpage

% =========================================================================
\section{CHƯƠNG 1: GIỚI THIỆU VỀ KINH TẾ HỌC}
% =========================================================================

\subsection{Lý thuyết cốt lõi của Bài giảng}
* **Sự khan hiếm:** Xảy ra khi mong muốn vô hạn của con người vượt quá nguồn lực hữu hạn. Sự khan hiếm buộc con người phải đưa ra sự lựa chọn, mỗi sự lựa chọn đều phải đánh đổi và phát sinh \textbf{chi phí cơ hội}.
* **Chi phí cơ hội (Opportunity Cost):** Là giá trị của phương án thay thế tốt nhất bị bỏ qua khi đưa ra sự lựa chọn.
  * \textit{Ví dụ trong thực tế}: Đi học Đại học 1 năm tiêu tốn học phí, giáo trình và cả \textbf{tiền lương bị bỏ lỡ} do không đi làm. Xem phim tiêu tốn tiền vé và \textbf{giá trị thời gian} trong rạp.
* **Phát biểu Thực chứng (Positive) và Chuẩn tắc (Normative):**
  \begin{itemize}
      \item \textbf{Thực chứng}: Mô tả thế giới khách quan, trả lời câu hỏi "Là gì?". Có thể kiểm chứng bằng số liệu. 
      \\ \textit{Ví dụ}: "Lạm phát tăng khi chính phủ in quá nhiều tiền."
      \item \textbf{Chuẩn tắc}: Đưa ra khuyến nghị hoặc đánh giá chủ quan, trả lời câu hỏi "Nên như thế nào?". 
      \\ \textit{Ví dụ}: "Chính phủ nên cắt giảm tốc độ lạm phát." hoặc "Nên miễn học phí cho tất cả các cấp."
  \end{itemize}

\subsection{Bài tập mẫu từ Slide Bài giảng}

\paragraph{Ví dụ 1 (Đường PPF và Chi phí cơ hội tăng dần):}
Cho bảng phối hợp sản lượng tối đa giữa 2 sản phẩm $X$ và $Y$ của một nền kinh tế:
\begin{center}
\begin{tabular}{cccccc}
\toprule
\textbf{Phối hợp} & \textbf{A} & \textbf{B} & \textbf{C} & \textbf{D} & \textbf{E} \\ \midrule
\textbf{X} & 0 & 50 & 100 & 150 & 200 \\
\textbf{Y} & 100 & 90 & 75 & 50 & 0 \\
\bottomrule
\end{tabular}
\end{center}
Hãy tính chi phí cơ hội biên để sản xuất thêm sản phẩm $X$ qua từng giai đoạn và nhận xét.

\paragraph{Lời giải:}
Chi phí cơ hội để sản xuất thêm $1$ đơn vị sản phẩm $X$ ($OC_X$) được tính theo công thức:
$$OC_X = \left| \frac{\Delta Y}{\Delta X} \right|$$
\begin{itemize}
    \item \textbf{Từ A sang B:} $OC_X = \frac{100 - 90}{50 - 0} = \frac{10}{50} = 0.2$ đơn vị Y.
    \item \textbf{Từ B sang C:} $OC_X = \frac{90 - 75}{100 - 50} = \frac{15}{50} = 0.3$ đơn vị Y.
    \item \textbf{Từ C sang D:} $OC_X = \frac{75 - 50}{150 - 100} = \frac{25}{50} = 0.5$ đơn vị Y.
    \item \textbf{Từ D sang E:} $OC_X = \frac{50 - 0}{200 - 150} = \frac{50}{50} = 1.0$ đơn vị Y.
\end{itemize}
\textbf{Nhận xét:} Chi phí cơ hội của $X$ tăng dần ($0.2 \rightarrow 0.3 \rightarrow 0.5 \rightarrow 1.0$) biểu thị đường PPF có dạng \textbf{cong lồi hướng ra ngoài} gốc tọa độ do nguồn lực sản xuất không thể chuyển đổi hoàn hảo giữa hai ngành.

\newpage

% =========================================================================
\section{CHƯƠNG 2: CUNG, CẦU VÀ CO GIÃN}
% =========================================================================

\subsection{Các công thức co giãn quan trọng}
\begin{itemize}
    \item \textbf{Co giãn khoảng của cầu theo giá:}
    $$E_D = \frac{Q_2 - Q_1}{P_2 - P_1} \times \frac{P_1 + P_2}{Q_1 + Q_2}$$
    \item \textbf{Co giãn thu nhập ($E_I$):} Phân loại hàng hóa:
    \begin{itemize}
        \item $E_I < 0$: Hàng cấp thấp (thứ cấp).
        \item $0 < E_I < 1$: Hàng thiết yếu.
        \item $E_I > 1$: Hàng cao cấp (xa xỉ).
    \end{itemize}
    \item \textbf{Co giãn chéo ($E_{XY}$):} Phân loại mối quan hệ:
    \begin{itemize}
        \item $E_{XY} > 0$: Hai hàng hóa thay thế.
        \item $E_{XY} < 0$: Hai hàng hóa bổ sung.
        \item $E_{XY} = 0$: Hai hàng hóa độc lập.
    \end{itemize}
\end{itemize}

\subsection{Bài tập mẫu từ Slide Bài giảng}

\paragraph{Ví dụ 1 (Tính co giãn khoảng từ số liệu):}
Cho bảng số liệu nhu cầu của người tiêu dùng:
\begin{center}
\begin{tabular}{ccc}
\toprule
\textbf{Trạng thái} & \textbf{P} & \textbf{$Q_D$} \\ \midrule
1 & 220 & 1200 \\
2 & 180 & 1600 \\
\bottomrule
\end{tabular}
\end{center}
Tính độ co giãn khoảng của cầu theo giá.

\paragraph{Lời giải:}
Áp dụng công thức co giãn trung điểm:
$$E_D = \frac{1600 - 1200}{180 - 220} \times \frac{220 + 180}{1200 + 1600} = \frac{400}{-40} \times \frac{400}{2800} = -10 \times \frac{1}{7} \approx -1.43$$
Vì $|E_D| = 1.43 > 1$, cầu co giãn nhiều.

\paragraph{Ví dụ 2 (Đánh giá Độ co giãn thu nhập):}
\begin{itemize}
    \item \textbf{Tình huống A}: Khi thu nhập giảm $25\%$, lượng mua hàng hóa $A$ giảm $10\%$.
    \\ $\Rightarrow E_I = \frac{-10\%}{-25\%} = 0.4 \implies 0 < E_I < 1 \Rightarrow A$ là \textbf{hàng thiết yếu} (hàng thông thường).
    \item \textbf{Tình huống B}: Khi thu nhập tăng $18\%$, lượng mua hàng hóa $B$ giảm $5\%$.
    \\ $\Rightarrow E_I = \frac{-5\%}{18\%} \approx -0.28 \implies E_I < 0 \Rightarrow B$ là \textbf{hàng thứ cấp}.
    \item \textbf{Tình huống C}: Khi thu nhập tăng $25\%$, lượng mua hàng hóa $C$ tăng $30\%$.
    \\ $\Rightarrow E_I = \frac{30\%}{25\%} = 1.2 \implies E_I > 1 \Rightarrow C$ là \textbf{hàng xa xỉ} (hàng cao cấp).
\end{itemize}

\paragraph{Ví dụ 3 (Tính toán tổng hợp co giãn và Cân bằng thị trường):}
Cho hàm cầu và cung thị trường: $Q_D = -0.5P + 200$ và $Q_S = 1.5P - 10$.
\begin{enumerate}
    \item Xác định giá và lượng cân bằng tự do.
    \item Tính độ co giãn của cung và cầu tại điểm cân bằng.
    \item Để tăng doanh thu, nhà sản xuất nên tăng hay giảm giá?
\end{enumerate}
\paragraph{Lời giải:}
\begin{enumerate}
    \item \textbf{Điểm cân bằng}: $Q_D = Q_S \implies -0.5P + 200 = 1.5P - 10 \implies 2P = 210 \implies P_0 = 105$.
    \\ $\implies Q_0 = 1.5(105) - 10 = 147.5$.
    \item \textbf{Độ co giãn tại điểm cân bằng}:
    \begin{itemize}
        \item Cầu: $E_D = Q'_D(P) \times \frac{P_0}{Q_0} = -0.5 \times \frac{105}{147.5} \approx -0.356$.
        \item Cung: $E_S = Q'_S(P) \times \frac{P_0}{Q_0} = 1.5 \times \frac{105}{147.5} \approx 1.068$.
    \end{itemize}
    \item \textbf{Phân tích doanh thu}: Vì $|E_D| \approx 0.356 < 1$ (cầu co giãn ít), mối quan hệ giữa giá và doanh thu ($TR$) là đồng biến. Do đó, để tăng doanh thu, nhà sản xuất nên \textbf{tăng giá}.
\end{enumerate}

\newpage

% =========================================================================
\section{CHƯƠNG 3: LÝ THUYẾT HÀNH VI NGƯỜI TIÊU DÙNG}
% =========================================================================

\subsection{Phương pháp tính hữu dụng biên (MU) rời rạc}
Trong bài tập rời rạc, hữu dụng biên được tính dựa trên phần chênh lệch:
$$MU_i = \frac{\Delta U_i}{\Delta i} = \frac{U_{i2} - U_{i1}}{i_2 - i_1}$$

\subsection{Bài tập mẫu từ Slide Bài giảng}

\paragraph{Ví dụ 1 (Tính MU rời rạc có khoảng cách đơn vị):}
Cho tổng hữu dụng của $X$: $U_X = [9, 16, 21, 24, 25]$ tương ứng với lượng $x = [1, 2, 3, 4, 5]$.
\begin{itemize}
    \item $MU_X(1) = 9$
    \item $MU_X(2) = (16-9)/(2-1) = 7$
    \item $MU_X(3) = (21-16)/(3-2) = 5$
    \item $MU_X(4) = (24-21)/(4-3) = 3$
    \item $MU_X(5) = (25-24)/(5-4) = 1$
\end{itemize}
Giá trị hữu dụng biên giảm dần ($9 \to 7 \to 5 \to 3 \to 1$) minh họa cho quy luật hữu dụng biên giảm dần.

\paragraph{Ví dụ 2 (Tính MU rời rạc có khoảng cách bước bất kỳ):}
Cho bảng tổng hữu dụng của $y$ như sau. Hãy điền các giá trị $MU_y$ còn thiếu:
\begin{center}
\begin{tabular}{cccccc}
\toprule
\textbf{$y$} & 0 & 2 & 5 & 8 & 10 \\ \midrule
\textbf{$U_y$} & 0 & 20 & 44 & 62 & 70 \\
\textbf{$MU_y$} & 0 & ? & ? & ? & ? \\
\bottomrule
\end{tabular}
\end{center}
\paragraph{Lời giải:}
Ta tính lượng tăng thêm trên mỗi đơn vị đầu vào $\Delta U / \Delta y$:
\begin{itemize}
    \item Tại $y = 2$: $MU_y = \frac{20 - 0}{2 - 0} = 10$.
    \item Tại $y = 5$: $MU_y = \frac{44 - 20}{5 - 2} = \frac{24}{3} = 8$.
    \item Tại $y = 8$: $MU_y = \frac{62 - 44}{8 - 5} = \frac{18}{3} = 6$.
    \item Tại $y = 10$: $MU_y = \frac{70 - 62}{10 - 8} = \frac{8}{2} = 4$.
\end{itemize}

\paragraph{Ví dụ 3 (Tìm ngược hàm tổng thỏa dụng rời rạc):}
Cho bảng hữu dụng biên của $y$ như sau. Điền các giá trị $U_y$ còn thiếu:
\begin{center}
\begin{tabular}{cccccc}
\toprule
\textbf{$y$} & 0 & 1 & 3 & 5 & 6 \\ \midrule
\textbf{$MU_y$} & 0 & 9 & 7 & 5 & 3 \\
\textbf{$U_y$} & 0 & ? & ? & ? & ? \\
\bottomrule
\end{tabular}
\end{center}
\paragraph{Lời giải:}
Áp dụng công thức $U_{new} = U_{old} + MU \times \Delta y$:
\begin{itemize}
    \item Tại $y = 1$: $U_y = 0 + 9 \times 1 = 9$.
    \item Tại $y = 3$: $U_y = 9 + 7 \times (3 - 1) = 9 + 14 = 23$.
    \item Tại $y = 5$: $U_y = 23 + 5 \times (5 - 3) = 23 + 10 = 33$.
    \item Tại $y = 6$: $U_y = 33 + 3 \times (6 - 5) = 33 + 3 = 36$.
\end{itemize}

\paragraph{Ví dụ 4 (Tìm phối hợp tiêu dùng tối ưu bậc cao):}
Một người tiêu dùng có tổng hữu dụng $TU = (4X-8)Y$. Người này có thu nhập hàng tháng $I = 30$ triệu đồng dùng mua $X$ và $Y$ với đơn giá $P_X = 3$ triệu đồng/đơn vị và $P_Y = 6$ triệu đồng/đơn vị. Tìm rổ hàng tối ưu để thỏa dụng tối đa.
\paragraph{Lời giải:}
\begin{enumerate}
    \item Tính hữu dụng biên bằng đạo hàm riêng:
    $$MU_X = \frac{\partial TU}{\partial X} = 4Y$$
    $$MU_Y = \frac{\partial TU}{\partial Y} = 4X - 8$$
    \item Lập điều kiện cân bằng tiêu dùng:
    $$\frac{MU_X}{P_X} = \frac{MU_Y}{P_Y} \implies \frac{4Y}{3} = \frac{4X-8}{6} \implies 8Y = 4X - 8 \implies X = 2Y + 2$$
    \item Lập phương trình ngân sách:
    $$3X + 6Y = 30 \implies X + 2Y = 10$$
    \item Giải hệ phương trình thế (1) vào (2):
    $$(2Y + 2) + 2Y = 10 \implies 4Y = 8 \implies Y^* = 2$$
    $$\implies X^* = 2(2) + 2 = 6$$
    \item \textbf{Kết quả:} Rổ hàng tối ưu là $(X, Y) = (6, 2)$.
\end{enumerate}

\newpage

% =========================================================================
\section{CHƯƠNG 4: LÝ THUYẾT SẢN XUẤT VÀ CHI PHÍ}
% =========================================================================

\subsection{Mối quan hệ cận biên và trung bình}
\paragraph{Quy luật quan hệ biên - trung bình:}
Đại lượng biên lớn hơn đại lượng trung bình sẽ kéo đại lượng trung bình đi lên; và ngược lại.
\begin{itemize}
    \item Nếu $MP_L > AP_L \implies AP_L$ đang tăng.
    \item Nếu $MP_L < AP_L \implies AP_L$ đang giảm.
    \item Nếu $MP_L = AP_L \implies AP_L$ đạt cực đại.
\end{itemize}
\paragraph{Phép ẩn dụ tuổi trung bình của lớp học:}
Lớp có 40 sinh viên đều 18 tuổi (tuổi trung bình là 18).
\begin{itemize}
    \item Nếu thêm 1 sinh viên mới có tuổi biên là 19 ($19 > 18$) $\implies$ Tuổi trung bình lớp tăng.
    \item Nếu thêm 1 sinh viên mới có tuổi biên là 16 ($16 < 18$) $\implies$ Tuổi trung bình lớp giảm.
\end{itemize}

\subsection{Bài tập mẫu từ Slide Bài giảng}

\paragraph{Ví dụ 1 (Điền bảng năng suất ngắn hạn ngắn):}
Hãy hoàn thành các ô trống trong bảng năng suất ngắn hạn sau (với $K=3$ cố định):
\begin{center}
\begin{tabular}{ccccc}
\toprule
\textbf{K} & \textbf{L} & \textbf{Q = TP} & \textbf{$AP_L$} & \textbf{$MP_L$} \\ \midrule
3 & 0 & 0 & - & - \\
3 & 5 & 200 & ? & ? \\
3 & 10 & ? & 35 & ? \\
3 & 15 & ? & ? & 40 \\
\bottomrule
\end{tabular}
\end{center}
\paragraph{Lời giải:}
\begin{itemize}
    \item Tại $L = 5$: $AP_L = 200/5 = 40$. $MP_L = (200-0)/(5-0) = 40$.
    \item Tại $L = 10$: $AP_L = 35 \implies Q = 35 \times 10 = 350$. $MP_L = (350-200)/(10-5) = 150/5 = 30$.
    \item Tại $L = 15$: $MP_L = 40 \implies \Delta Q = 40 \times 5 = 200 \implies Q = 350 + 200 = 550$.
    \\ $AP_L = 550/15 \approx 36.67$.
\end{itemize}

\paragraph{Ví dụ 2 (Điền bảng chi phí ngắn hạn hoàn chỉnh):}
Hãy điền đầy đủ các ô trống của bảng số liệu sau (đáp án từ Ví dụ 6 bài giảng):
\begin{center}
\begin{tabular}{cccccccc}
\toprule
\textbf{Q} & \textbf{TC} & \textbf{TFC} & \textbf{TVC} & \textbf{AC (ATC)} & \textbf{AVC} & \textbf{AFC} & \textbf{MC} \\ \midrule
\textbf{0} & 100 & 100 & 0 & - & - & - & - \\
\textbf{10} & 200 & 100 & 100 & 20 & 10 & 10 & 10 \\
\textbf{20} & 280 & 100 & 180 & 14 & 9 & 5 & 8 \\
\textbf{30} & 450 & 100 & 350 & 15 & 11.6 & 3.4 & 17 \\
\textbf{40} & 660 & 100 & 560 & 16.5 & 14 & 2.5 & 21 \\
\textbf{50} & ? & 100 & ? & ? & ? & ? & 25 \\
\bottomrule
\end{tabular}
\end{center}
\paragraph{Lời giải:}
Ta điền dòng cuối cùng tại $Q = 50$:
\begin{itemize}
    \item Từ $Q=40$ lên $Q=50$, với $MC = 25$:
    $$\Delta TC = MC \times \Delta Q = 25 \times 10 = 250 \implies TC_{50} = 660 + 250 = 910$$
    \item $TVC_{50} = TC_{50} - TFC = 910 - 100 = 810$.
    \item $AC_{50} = 910 / 50 = 18.2$.
    \item $AVC_{50} = 810 / 50 = 16.2$.
    \item $AFC_{50} = 100 / 50 = 2.0$.
\end{itemize}

\paragraph{Ví dụ 3 (Bài toán hàm chi phí liên tục):}
Cho hàm số tổng chi phí ngắn hạn: $TC = 2Q^3 - 20Q^2 + 100Q + 2500$.
Hãy xác định các hàm chi phí: $TFC$, $TVC$, $AFC$, $AVC$, $ATC$, $MC$.
\paragraph{Lời giải:}
\begin{itemize}
    \item $TFC = 2500$ (hệ số tự do).
    \item $TVC = 2Q^3 - 20Q^2 + 100Q$.
    \item $AFC = \frac{TFC}{Q} = \frac{2500}{Q}$.
    \item $AVC = \frac{TVC}{Q} = 2Q^2 - 20Q + 100$.
    \item $ATC = \frac{TC}{Q} = 2Q^2 - 20Q + 100 + \frac{2500}{Q}$.
    \item $MC = (TC)' = 6Q^2 - 40Q + 100$.
\end{itemize}

\end{document}